The proofs below repeatedly rely on two key ideas: Taylor series
expansions, and uniform asymptotic control on the terms that arise.
In this section, we introduce these ideas and prove two key
lemmas.  The first, \lemref{taylor_residual}, gives the form
of a Taylor series expansion.  The second, \lemref{ulln},
controls the asymptotic behavior of the log likelihood
and its derivatives in a neighborhood of $\thetatrue$.


%%%%%%%%%%%%%%%%%%%%%%%%%%%%%%%%%%%%%%%%%
%%%%%%%%%%%%%%%%%%%%%%%%%%%%%%%%%%%%%%%%%
%%%%%%%%%%%%%%%%%%%%%%%%%%%%%%%%%%%%%%%%%


\subsubsection{Taylor series expansions}

Our analysis is based strongly on Taylor series expansions.
We will use the following two minor variants of standard analysis tools.
The first is a multivariate version of the integral form of the Taylor series
residual for a vector-valued function.

\begin{defn}\deflabel{taylor_residual}
%
For a $k$-times differentiable function $\phi(\theta)$, a fixed $\thetahat$, and
and $\theta$, define the residual functional
%
\begin{align*}
%
\tresid{k}{\phi, \thetahat, \theta} :=
    \frac{1}{(k-1)!} \int_0^1 (1 - t)^{k-1}
        \phigrad{k}(\thetahat + t (\theta - \thetahat)) dt.
%
\end{align*}
%
\end{defn}
%
The residual $\tresid{k}{\phi, \thetahat, \theta}$ is an array of the same
dimension as the $k$-th order derivative, $\phigrad{k}(\theta)$.  Observe that,
for any norm $\norm{\cdot}$,
%
\begin{align*}
    %
    \norm{\tresid{k}{\phi, \thetahat, \theta}} \le
        \frac{1}{(k-1)!} \sup_{\thetatil: \norm{\thetatil - \thetahat}_\infty \le
                                          \norm{\theta - \thetahat}_\infty}
                            \norm{\phigrad{k}(\thetatil)}.
    %
\end{align*}
%
\begin{lem}\lemlabel{taylor_residual}
%
A $K$-times continuously differentiable function $\phi(\theta)$ has Taylor
series residual
%
\begin{align*}
\phi(\theta) =
    \sum_{k=0}^{K-1} \frac{1}{k!} \phigrad{k}(\thetahat) (\theta - \thetahat)^k +
    \tresid{K}{\phi, \thetahat, \theta} (\theta - \thetahat)^K.
\end{align*}
%
\begin{proof}
%
Let $\phi_d(\theta)$ denote the $d$-th component of the vector $\phi(\theta)$
(contrast this notation with for the $d$-th order derivative
$\phigrad{d}(\theta)$). Consider the scalar to scalar map $t \mapsto
\phi_d(\thetahat + t (\theta - \thetahat))$. Taking the integral remainder form
of the $K - 1$-th order Taylor series expansion of this scalar map
\cite[Appendix B.2]{dudley:2018:analysis} and stacking the result into a vector
gives the desired result.
%
\end{proof}
%
\end{lem}



%%%%%%%%%%%%%%%%%%%%%%%%%%%%%%%%%%%%%%%%%
%%%%%%%%%%%%%%%%%%%%%%%%%%%%%%%%%%%%%%%%%
%%%%%%%%%%%%%%%%%%%%%%%%%%%%%%%%%%%%%%%%%

\subsubsection{Log likelihood asymptotics}

To use \lemref{taylor_residual} to study $\likhat(\theta)$, we will need to
expand \defref{loglik_def} with the following quantities.

%%%%%%%%%%%%%%%%%%%%%%%%%%%%%%%%%%%%%%%%
%%%%%%%%%%%%%%%%%%%%%%%%%%%%%%%%%%%%%%%%
\begin{defn}\deflabel{loglik_details_def}
\begin{align*}
%
\likhatk{k}(\theta) :={}&
   \frac{1}{N} \sumn \ellgrad{k}(\x_n \vert \theta) +
    %\frac{1}{N} \frac{\prior_{(k)}(\theta)}{\prior(\theta)} % NOT CORRECT
    \frac{1}{N} \logpriorgrad{k}(\theta)
&
\likk{k}(\theta) :={}& \int \ellgrad{k}(\xn \vert \theta) d\fdist(\xn)
\\
\infoev :={}& \textrm{The minimum eigenvalue of }\info &
\infoevhat :={}& \textrm{The minimum eigenvalue of }\infohat
%
\end{align*}
%
Note that, following our convention, $\likhatk{k}(\theta)$ is the $k$--th
partial derivative of $\likhat(\theta)$ with respect to $\theta$.

The term $\likk{k}(\theta)$ is a potential abuse of notation, since it is only
the partial derivative of $\lik(\theta)$ under the assumption that we can
interchange integration and differentiation.  Though we will in fact be assuming
sufficient conditions for this interchange (\assuref{bayes_clt}
\itemref{loglik_ulln}), our proofs will simply use the above definition of
$\likk{k}(\theta)$ directly.
\end{defn}
%
Using \lemref{taylor_residual}, we can then form a Taylor series expansion
around $\thetahat$, an expression that we will analyze carefully in
\appref{bayes_clt_proof}:
%
\begin{align}
\likhat(\theta) ={}&
    \likhat(\thetahat) +
    \likhatk{1}(\thetahat)(\theta - \thetahat) +
    \frac{1}{2}\likhatk{2}(\thetahat)(\theta - \thetahat)^2 +
\nonumber
\\&
    \frac{1}{6}\likhatk{3}(\thetahat)(\theta - \thetahat)^3 +
    \tresid{4}{\likhat, \thetahat, \theta} (\theta - \thetahat)^4.
    \eqlabel{loglik_series_example}
\end{align}
%

We now turn to the task of controlling the kinds of terms that arise in
\eqref{loglik_series_example}. One of our key tools in the proof of
\lemref{ulln} will be \emph{uniform laws of large numbers} (ULLNs), a concept
which we make concrete in \defref{ulln}. Typically, ULLNs require that the
sample average $\meann \phi(\x_n, \theta)$ converge to
$\expect{\fdist(\xn)}{\phi(\x_n, \theta)}$ uniformly over $\theta$ in some
compact set.
%
%%%%%%%%%%%%%%%%%%%%%%%%%%%%%%%%%%%%%%%%%%%%%
%%%%%%%%%%%%%%%%%%%%%%%%%%%%%%%%%%%%%%%%%%%%%
\begin{defn}\deflabel{ulln}
%
Recall the definition of $\thetaball{\delta}$ from \assuref{bayes_clt}. We say a
``uniform law of large numbers'', or ``ULLN'', holds for a function
$\phi(\theta, \xn)$ if, for some $\delta > 0$,
%
\begin{align*}
%
\sup_{\theta \in \thetaball{\delta}}
\norm{\meann \phi(\theta, \x_n) -
\expect{\fdist(\xn)}{\phi(\theta, \xn)}}_2
\plim 0,
%
\end{align*}
%
where $\expect{\fdist(\xn)}{\phi(\theta, \xn)}$ is finite.
%
\end{defn}
%%%%%%%%%%%%%%%%%%%%%%%%%%%%%%%%%%%%%%%%%%%%%

In general, a ULLN requires that the set of functions $\left\{\theta \mapsto
\phi(\theta, \xvec): \norm{\theta - \thetatrue}_2 < \delta \right\}$ be a
``Glivenko-Cantelli class'' with finite covering number (Chapter 19 of
\citet{van:2000:asymptotic}, \citet{vaart:2013:empiricalprocesses}). A simple
condition for a uniform law of large numbers to hold is that $\phi(\theta, \xn)$
be dominated by an absolutely integrable function of $\xn$ alone in a compact
neighborhood of $\thetatrue$ (Example 19.7 of \citet{van:2000:asymptotic}).

We now state a key lemma controlling the asymptotic behavior of the log
likelihood and its derivatives.

%%%%%%%%%%%%%%%%%%%%%%%%%%%%%%%%%%%%%%%%%%%%%%%%%%%
%%%%%%%%%%%%%%%%%%%%%%%%%%%%%%%%%%%%%%%%%%%%%%%%%%%
%%%%%%%%%%%%%%%%%%%%%%%%%%%%%%%%%%%%%%%%%%%%%%%%%%%

\begin{lem}\lemlabel{ulln}
%
Suppose that there exists $\deltaulln > 0$, $M(\xn)$, and $K \ge 1$ such that
%
\begin{align}
%
\sup_{\theta \in \thetaball{\deltaulln}}
    \norm{\ellgrad{K}(\xn \vert \theta)}_2 \le M(\xn)
\quad\textrm{with}\quad
\expect{\fdist(\xn)}{M(\xn)^2} \le \infty.
\eqlabel{bounded_lipschitz}
%
\end{align}
%
Addionally assume that $\pi(\theta)$ is non-zero and $K-1$ times
continuously differentiable in $\thetaball{\deltaulln}$.
%
Then the following results hold.  First, for $k \in \{0, \ldots, K\}$,
%
\begin{align}
\sup_{\theta \in \thetaball{\deltaulln}}
    \meann \norm{\ellgrad{k}(\x_n \vert \theta)}_2^2
     &={} \ordp{1}
% \nonumber\\&\hspace{-3em}
     \eqlabel{means_op1}.
%
\end{align}
%
Then, for $k \in \{0, \ldots, K - 1\}$,
%
\begin{align}
%
\sup_{\theta \in \thetaball{\deltaulln}}
    \norm{\meann \left( \ellgrad{k}(\x_n \vert \theta) -
           \expect{\fdist(\xn)}{\ellgrad{k}(\xn \vert \theta)} \right)
           }_2 &\plim{} 0
         \eqlabel{glivenko_cantelli} \\
%
\sup_{\theta \in \thetaball{\deltaulln}}
    \norm{\frac{1}{\sqrt{N}} \sumn \left( \ellgrad{k}(\x_n \vert \theta) -
           \expect{\fdist(\xn)}{\ellgrad{k}(\xn \vert \theta)}\right)
           }_2 &={} \ordp{1}
           \eqlabel{donsker} \\
%
\sup_{\theta \in \thetaball{\deltaulln}}
    \norm{\likhatk{k}(\theta) - \likk{k}(\theta)}_2 &\plim{} 0.
         \eqlabel{ulln}
%
\end{align}
%
%%%%%%%%%%%%%%%%%%%%%%%%%%%%%%%%%%%%%%%%%%
%%%%%%%%%%%%%%%%%%%%%%%%%%%%%%%%%%%%%%%%%%
\begin{proof}
%
First, note that
%
\begin{align*}
%
\sup_{\theta \in \thetaball{\deltaulln}}
    \meann \norm{\ellgrad{K}(\x_n \vert \theta)}_2^2 \le{}&
\meann \sup_{\theta \in \thetaball{\deltaulln}}
    \norm{\ellgrad{K}(\x_n \vert \theta)}_2^2
\quad\textrm{and}\\
\expect{\fdist(\xn)}{
    \sup_{\theta \in \thetaball{\deltaulln}}
        \norm{\ellgrad{K}(\xn \vert \theta)}^2_2
} \le{}& \expect{\fdist(\xn)}{M(\xn)^2} < \infty,
%
\end{align*}
%
so \eqref{means_op1} holds for $k = K$ by the strong law of large numbers.

We now prove the remaining relations inductively. For a scalar--valued,
differentiable function $f(\theta)$, and any two $\theta, \theta'$ in
$\thetaball{\deltaulln}$, we can write
%
\begin{align*}
%
\abs{f(\theta) - f(\theta')} ={}&
\abs{
    \int_0^1 \fracat{\partial f(\theta' + \t (\theta - \theta'))}
                    {\partial \t}{\t} d\t
}
\\={}&
\abs{
    \int_0^1 f_{(1)}(\theta' + \t (\theta - \theta')) (\theta - \theta') d\t
}
\\\le{}&
    \norm{\theta - \theta'}_2
    \sup_{\thetatil \in \thetaball{\deltaulln}}
        \norm{f_{(1)}(\thetatil)}_2.
%
\end{align*}
%
Applying the previous formula componentwise to $\ellgrad{K-1}(\xn \vert
\theta)$, and loosening the right hand side to contain all of
$\norm{\ellgrad{K}(\xn \vert \thetatil)}_2$ rather than only a particular
component, gives
%
\begin{align}
%
\norm{\ellgrad{K-1}(\xn \vert \theta)}_\infty \le{}&
    \norm{\theta - \theta'}_2
    \sup_{\thetatil \in \thetaball{\deltaulln}}
        \norm{\ellgrad{K}(\xn \vert \thetatil)}_2
\nonumber \\\le{}&
\norm{\theta - \theta'}_2 M(\xn). \eqlabel{ellgrad_linf_bound}
%
\end{align}
%
It follows that each component of $\ellgrad{K-1}(\xn \vert \theta)$ is both
Glivenko-Cantelli and Donsker, and so \eqref{glivenko_cantelli,donsker} hold for
$k = K - 1$ (see \cite[Example 19.7]{van:2000:asymptotic},
as well as the discussion in \defref{ulln} above).  Further,
%
\begin{align*}
%
\sup_{\theta \in \thetaball{\deltaulln}}
    \norm{\ellgrad{K-1}(\xn \vert \theta)}_2
\le{}&
\sup_{\theta \in \thetaball{\deltaulln}}
    \sqrt{\thetadim} \norm{\ellgrad{K-1}(\xn \vert \theta)}_\infty
\le{}
    2 \sqrt{\thetadim}  \deltaulln M(\xn),
%
\end{align*}
%
where the final line applies \eqref{ellgrad_linf_bound} and
$
\sup_{\theta, \theta' \in \thetaball{\deltaulln}} \norm{\theta - \theta'}_2 \le 2 \deltaulln.
$
It follows that \eqref{means_op1} holds for $k = K - 1$.

Since \eqref{bounded_lipschitz} holds for $k=K-1$, the remaining results follow
by proceeding inductively on $K-1$, $K-2$, and so on.

Finally, by continuity of $\prior(\theta)$
(\assuref{bayes_clt} \itemref{prior_smooth}) and compactness of
$\thetaball{\deltaulln}$, we have that, as $N \rightarrow \infty$,
$$
\frac{1}{N} \sup_{\theta \in \thetaball{\deltaulln}} \norm{\logpriorgrad{k}(\theta)}_2^2  \rightarrow 0.
$$
Then \eqref{ulln} follows from the preceding display, Cauchy-Schwarz, and \eqref{glivenko_cantelli}.
%
\end{proof}
%

\end{lem}
